\chapter[Cohomological properness of sheaves of sets...]{Cohomological properness\\ of sheaves of sets and\\ of sheaves of non-commutative groups}
\label{XIII}
% original p. 344

\begin{center}
{by Mme M.\ \textsc{Raynaud}\footnote{After unpublished notes of
A\ptbl Grothendieck.}}
\end{center}
\vspace*{1cm}

This expos\'e proposes to use \'etale cohomology in order to
generalize certain results of \Ref{IX} and \Ref{X}. It also shows
how one can extend to sheaves of not necessarily commutative groups
those results of~SGA~5~II which still make sense for such
sheaves. The notions of \'etale cohomology expounded in~SGA~4 are
assumed known.

The main result \eqref{XIII.2.4} gives an important example
of a non-proper morphism $f\colon~U\to~S$ which is ``cohomologically
proper in dimension~$\leq1$'', that is to say such that, for
certain sheaves of groups $F$ on~$U$ (in the sense of the \'etale
topology), the formation of~$f_*F$ and $\R^1f_*F$ commutes with every
change of base~$S'\to~S$. This property is indeed
satisfied by the open set $U$ of a scheme $X$ proper over~$S$,
complementary to a divisor $D$ with normal crossings
relative to~$S$, at least if one requires $F$ to be
constant finite, of order prime to the residue
characteristics of~$S$. If one no longer assumes $F$ of order prime
to the residue characteristics of~$S$, one has an
analogous result on replacing $\R^1f_*F$ by the subsheaf
$\R^1_{\tame}f_*F$ obtained by restricting oneself to considering the torsors
under~$F$ ``tamely ramified on~$X$ relative
to~$S$''. In particular this makes it possible to show that the tamely
ramified fundamental group of a proper and smooth algebraic
curve over a separably closed field, with a finite number of
closed points removed, is topologically of finite type
\eqref{XIII.2.12}.

% original p. 345
No.~\Ref{XIII.4} is devoted to the homotopy exact sequence and to
the K\"unneth formula.

Finally an appendix gives useful variants of Abhyankar's lemma
proved in X.\Ref{X.3.6}.

\setcounter{section}{-1}

\section{Recollections on the theory of stacks}
\label{XIII.0}

We shall use in what follows the theory of stacks
expounded in \cite{XIII.1} and \cite{XIII.2}. We restrict ourselves to the case of the
\'etale site of a scheme. Given a scheme $X$, let us denote by
$X_\et$ the \'etale site of~$X$. Recall that a
stack $F$
on~$X$ is a fibered category over~$X_\et$ such
that, for every scheme $X'$ \'etale over~$X$ and for every pair
of objects
$x$, $y$ of the fiber $F_{X'}$, the presheaf
$\SheafHom_{X'}(x,y)$ is a sheaf, and such that, for every
surjective \'etale morphism $X'' \to X'$, every object of~$F_{X''}$
endowed with a descent datum relative to $X'' \to X'$ is the
inverse image of an object of~$F_{X'}$.

We denote by $F(X')$ the category of cartesian sections of
$F/X'$. More generally, if $\Sch_X$ is the category
of schemes over~$X$ endowed with the \'etale topology, the
stack $F$ can be extended to a stack $\cal{F}$ on~$\Sch_X$ and,
for every morphism $f\colon X' \to X$, we again denote by $F(X')$ the
category of cartesian sections of this stack $\cal{F}$
over~$X'$.

A gerbe
is a stack such that, for every scheme $X'$ \'etale
over~$X$ and for every pair of objects $x$, $y$ of~$F_{X'}$, every
morphism from~$x$ to $y$ is an isomorphism and $x$ and $y$ are
locally isomorphic, and such that the set of objects $X'$ of
$X_\et$ such that $F_{X'}$ is nonempty is a refinement of
$X_\et$. For example the stack of torsors under a sheaf of groups
is a gerbe which, moreover, has a cartesian section.
Conversely a gerbe which has a section, \ie such that there
exists an object $x$ of~$F_X$, is equivalent to the stack of
torsors under the sheaf
% original p. 346
of groups $\SheafAut_X(x)$.

There is an obvious notion of subgerbe and of maximal subgerbe
of a stack~$F$. Given a cartesian section $x$ of
$F(X)$, there exists a unique maximal subgerbe $G_x$ of~$F$ such
that $x$ factors through~$G_x$. We call $G_x$ the
subgerbe generated by~$x$; it is by definition a trivial
gerbe.
The presheaf $SF$
defined by
$$
SF(X') = \{ \text{maximal subgerbes of }F|X' \}
$$
is a sheaf, called the sheaf of maximal subgerbes
of~$F$.
Let $O$ be the presheaf defined by
$$
O(X') = \{ \text{classes of objects of }F_{X'}\text{ mod. isomorphism}
\} \text{.}
$$
By associating to every object $x$ of~$F_{X'}$ the maximal subgerbe
of~$F|X'$ generated by~$x$, we obtain a morphism
$$
O \to SF \text{;}
$$
by \cite[III 2.1.4]{XIII.2}, this morphism makes~$SF$ a sheaf
associated to~$O$.

A stack $F$ is said to be \emph{constructible}
(\emph{\resp ind-$\LL$-finite},
$\LL$ being a set of prime numbers) if, for every
scheme $X'$ \'etale over~$X$ and for every object $x$ of~$F_{X'}$,
the same holds for the sheaf $\SheafAut_{X'}(x)$ \cite[VII 2.2.1]{XIII.2}. A stack is
said to be $1$-constructible
if it is constructible and if the sheaf of maximal subgerbes is
constructible.

\section{Cohomological properness}
\label{XIII.1}
%
\setcounter{subsection}{-1}
\subsection{}
\label{XIII.1.0}
%
Let $S$ be a scheme, $f \colon X \to Y$ a morphism of
$S$-schemes. If $S'$ is an $S$-scheme, we consider the
following diagram, all of whose squares are cartesian:
% original p. 347
\begin{equation*}
\label{eq:XIII.1.0.1}
\tag{\thesubsection.1}
\begin{array}{c}
\xymatrix{ X \ar[d]_f & \ar[l]_h X'
\ar[d]^{f'} \\ Y \ar[d] & \ar[l]_g Y' \ar[d] \\ S & \ar[l] \,S'. }
\end{array}
\end{equation*}

If $Y_1$ is a scheme \'etale over~$Y$, we set $X_1 =
X \times_Y Y_1$, $Y_1' = Y' \times_Y Y_1$, and we consider the
cartesian square
\begin{equation*}
\label{eq:XIII.1.0.2}
\tag{\thesubsection.2}
\begin{array}{c}
\xymatrix{ X_1 \ar[d]_{f_1} & \ar[l]_{h_1}
X_1' \ar[d]^{f_1'} \\ Y_1 & \ar[l]^{g_1} \,Y_1'. }
\end{array}
\end{equation*}

\begin{definition}
\label{XIII.1.1}
Let $F$ be a stack on~$X$. We say that $(F,f)$ is
\emph{cohomologically proper relative to~$S$ in dimension
$\leqslant -1$}
(\resp in dimension $\leqslant 0$, \resp in dimension $\leqslant 1$)
if, for every $S$-scheme $S'$, the canonical functor (defined
in the obvious way by the universal property of
the inverse image of stacks):
$$
g^* f_* F \to f'_* h^* F \qquad\text{ (\cf 1.0.1)}
$$
is faithful (\resp fully faithful, \resp an equivalence
of categories).
\end{definition}

If there is no possible confusion about~$S$, in particular if $S=Y$,
we say cohomologically proper instead of cohomologically proper
relative to~$S$.

\subsection{}
\label{XIII.1.2}
Let $F$ be a sheaf of sets on~$X$; let $\Phi$ be the stack in
discrete categories associated to $F$, \ie the stack whose
fiber over every scheme $X_1$ \'etale over~$X$ is the
discrete category having $F(X_1)$ as its set of objects. We
say that $(F,f)$ is cohomologically proper relative to~$S$ in
dimension $\leqslant -1$ (\resp in dimension $\leqslant 0$) if $(\Phi,
f)$ is cohomologically proper relative to~$S$ in dimension
$\leqslant 0$ (\resp in dimension $\leqslant 1$).

The
% original p. 348
canonical morphism
\begin{equation*}
\label{eq:XIII.1.2.1}
\tag{\thesubsection.1} { g^* f_* F \to f'_* h^* F }
\end{equation*}
gives, on passing to the associated stacks in discrete
categories, the canonical morphism
$$
g^* f_* \Phi \to f'_* h^* \Phi \text{.}
$$
Consequently, to say that $(F,f)$ is cohomologically proper relative
to~$S$ in dimension $\leqslant -1$ (\resp in dimension $\leqslant
0$) is equivalent to saying that, for every $S$-scheme $S'$, the
morphism \eqref{eq:XIII.1.2.1} is injective (\resp bijective).

\subsection{}
\label{XIII.1.3}
Let $F$ be a sheaf of groups on~$X$ and $\Phi$ the stack of
torsors on~$X$ with group $F$ \cite[II 2.3.2]{XIII.1}. We say that $(F,f)$ is
cohomologically proper relative to~$S$ in dimension
$\leqslant -1$ (\resp $\leqslant 0$, \resp $\leqslant 1$) if
$(\Phi,f)$ is cohomologically proper relative to~$S$ in
dimension $\leqslant -1$ (\resp $\leqslant 0$, \resp $\leqslant
1$). The condition of cohomological properness can be made explicit
as follows.

\begin{subproposition}
\label{XIII.1.3.1}
The notation is that of \eqref{eq:XIII.1.0.1} and
\eqref{eq:XIII.1.0.2}. Let $F$ be a sheaf of groups on~$X$. We
denote by $F'$ (\resp $F_1$, \resp $F_1'$, etc.) the inverse
image of~$F$ on~$X'$ (\resp on~$X_1$, \resp on~$X_1'$,
etc.). Then the following conditions are equivalent:
\begin{enumerate}
\item[(i)] $(F,f)$ is cohomologically proper relative to~$S$
in dimension $\leqslant -1$ (\resp $\leqslant 0$, \resp $\leqslant
1$).
\item[(ii)] For every morphism $S' \to S$, for every scheme $Y_1$
\'etale over~$Y$, and for every torsor $P$ on~$X_1$ with
group $F_1$, if $\leftexp{P\mkern-4mu}{F}_1$ denotes the group $F_1$
twisted by $P$ \textup{\cite[II 4.1.2.3]{XIII.1}}, the canonical morphism
$$
a_0 \colon g_1^*(f_{1 *}(\leftexp{P\mkern-4mu}{F}_1)) \to f'_{1
*}(\leftexp{P'\mkern-6mu}{F}'_1)
$$
is injective (\resp $a_0$ is bijective and the canonical morphism
$$
a_1 \colon g^*(\R^1 f_* F) \to \R^1 f'_* F'
$$
is
% original p. 349
injective, \resp $a_0$ and $a_1$ are bijective).
\item[(ii~bis)] For every morphism $S' \to S$, for every scheme
$Y_1$ \'etale over~$Y$, for every torsor $P$ on~$X_1$ with
group $F_1$, and for every torsor $R$ under $\leftexp{P\mkern-4mu}{F}_1$, the
canonical morphism
$$
\alpha_0 \colon g_1^*(f_{1 *}R) \to f'_{1 *} R'
$$
is injective (\resp $\alpha_0$ is bijective, \resp the morphisms
$\alpha_0$ and
$$
\alpha_1 \colon g_1^*(\R^1 f_{1 *}(\leftexp{P\mkern-4mu}{F}_1)) \to \R^1
f'_{1 *}(\leftexp{P'\mkern-6mu}{F}'_1)
$$
are bijective).
\end{enumerate}
\end{subproposition}

\subsubsection*{Proof}
(i)$\To$(ii~bis). By \cite[II 4.2.5]{XIII.1} every torsor $R$
with group $\leftexp{P\mkern-4mu}{F}_1$ is of the form $R = Q
\overset{F_1}{\wedge} P^\circ$, where $Q$ is a torsor with group
$F_1$ and $P^\circ$ the opposite of~$P$. We then have $R' \simeq Q'
\overset{F'_1}{\wedge} {P'}^\circ$. Let $\Phi$ be the stack of torsors
under $F$ and let $x$, $y$ (\resp~$x'$,~$y'$) be the objects of the
fiber category $(g^* f_* \Phi)_{Y_1'}$ (\resp $(f'_*
\Phi')_{Y_1'}$) associated to $P$, $Q$ (\resp $P'$,~$Q'$). We have
the relation
$$
Q\overset{F_1}{\wedge} P^\circ \simeq \SheafHom_{F_1}(P,Q) \text{,}
$$
and it follows that we have canonical isomorphisms
$$
\SheafHom_{Y'_1}(x,y) \simeq g_1^* f_{1 *}(Q
\overset{F_1}{\wedge} P^\circ) \text{,}\quad \SheafHom_{Y_1'}(x',y')
\simeq f'_{1 *}(Q' \overset{F'_1}{\wedge}{P'}^\circ)\text{.}
$$
Consequently the morphism $\alpha_0$ is identified with the morphism
$$
\SheafHom_{Y'_1}(x,y) \to \SheafHom_{Y'_1}(x',y') \text{,}
$$
from which it follows that, if $(F,f)$ is cohomologically proper
relative to~$S$ in dimension $\leqslant -1$, $\alpha_0$ is
injective and that, if $(F,f)$ is cohomologically proper in dimension
$\leqslant 0$, $\alpha_0$ is bijective.

Suppose now that $(F,f)$ is cohomologically proper
relative to~$S$ in dimension $\leqslant 1$, \ie that the
canonical morphism
% original p. 350
$$
\varphi \colon g^* f_* \Phi \to f'_* \Phi'
$$
is an equivalence. Let $G$ be the sheaf of maximal subgerbes
of the stack $f_* \Phi$ \hbox{\cite[III~2.1.8]{XIII.1}}; we then have an
isomorphism $G \simeq \R^1 f_* F$. Since $g^* G$ is the
sheaf of maximal subgerbes of~$g^* f_* \Phi$ \cite[III
2.1.5.5]{XIII.2}, the morphism $\alpha_1$ is obtained from
$\varphi|Y'_1$ by taking the sheaves of maximal subgerbes,
hence is an isomorphism.

(ii~bis)$\To$(ii). It suffices to show that, if the
morphisms $\alpha_0$ are bijective, then the morphisms $a_1$ are
injective. Let $Y'_1$ be a scheme \'etale over~$Y'$, and
$s$ and $t$ two elements of~$g^*(\R^1 f_* F)(Y'_1)$
having the same image in $\R^1 f'_* F'(Y'_1)$, and let us show that
we have $s=t$. The assertion is local for the \'etale topology of
$Y'_1$ and, taking into account the definition of the inverse image
$g^*(\R^1 f_* F)$, we may assume that $Y'_1$ is the inverse
image of a scheme $Y_1$ \'etale over~$Y$ and that $s$ and
$t$ come from torsors $P$ and $Q$ on~$X_1$. The hypothesis
made on~$s$ and $t$ then means that the inverse images $P'$ and
$Q'$ of~$P$ and $Q$ on~$X'_1$ are isomorphic locally for the
\'etale topology of~$Y'_1$. If we set $R =
\SheafHom_{F_1}(P,Q)$, the fact that the morphism
$$
g_1^* f_{1 *} R \to f'_{1 *} R'
$$
is bijective proves that $P$ and $Q$ are isomorphic locally for the
\'etale topology of~$Y_1$, hence that we have $s=t$.

(ii)$\To$(i). To prove that $\varphi$ is faithful
(\resp fully faithful), it suffices to show that, if $Y_1$ is
a scheme \'etale over~$Y$, if $P$, $Q$ are two torsors on~$X_1$ with group~$F_1$, and if $x$, $y$ (\resp $x'$, $y'$) are the objects
of~$(g^* f_* \Phi)_{Y'_1}$ (\resp $(f'_* \Phi')_{Y'_1}$)
associated to $P$, $Q$ (\resp $P'$, $Q'$), then the morphism
$$
a \colon \Hom(x,y) \to \Hom(x',y')
$$
is injective (\resp bijective). Now $a$ is identified with the
canonical morphism
$$
\H^0(Y'_1,g^*_1 f_{1 *} (Q \overset{F_1}{\wedge} P^\circ)) \to
\H^0(Y'_1, f'_{1 *}(Q' \overset{F'_1}{\wedge} {P'}^\circ))\text{.}
$$
If
% original p. 351
we have $\Hom(x,y) \neq \emptyset$, then $Q \overset{F_1}{\wedge}
P^\circ$ is a torsor under $\leftexp{P\mkern-4mu}{F}_1$ which is locally trivial over~$Y_1$; it follows that $f_{1 *}(Q \overset{F_1}{\wedge}
P^\circ)$ is a torsor under $f_{1*}(\leftexp{P\mkern-4mu}{F}_1)$, and that
$g^*_1 f_{1 *} (Q \overset{F_1}{\wedge} P^\circ)$ is a trivial
torsor. The morphism $a$ is then identified with the canonical morphism
$$
\H^0(Y_1, g^*_1 f_{1 *}(\leftexp{P\mkern-4mu}{F}_1)) \to \H^0 (Y'_1,
f'_{1 *}(\leftexp{P'\mkern-6mu}{F}'_1)) \text{.}
$$
The same holds if we have $\Hom(x',y') \neq \emptyset$ and if
$a_1$ is injective, for then $Q' \overset{F'_1}{\wedge} {P'}^\circ$
is trivial, and it follows from the injectivity of~$a_1$ that $P$
and $Q$ are locally isomorphic over~$Y_1$. We conclude that, if
$a_0$ is injective (\resp if $a_0$ is bijective and $a_1$ injective),
$\varphi$ is faithful (\resp fully faithful).

It remains to show that, if $a_0$ and $a_1$ are bijective, the
functor $\varphi$ is essentially surjective. Let $Y''$ be a
scheme \'etale over~$Y'$, $X'' = X' \times_{Y'} Y''$, and
let $P''$ be a torsor on~$X''$ with group $F'' = F'|X''$. We shall
show that there exists an element~$x$ of~$(g^* f_*
\Phi)_{Y''}$ whose image in $(f'_* \Phi')_{Y''}$ is isomorphic
to $P''$. Let $p''$ be the class of~$P''$. From the fact that $a_1$ is
surjective it follows that one can find a surjective \'etale
morphism $Y''_1 \to Y''$, an \'etale morphism $Y_1 \to Y$ such that
there is a morphism $Y''_1 \to Y'_1$, and a torsor $P_1$ on~$X_1$
with group $F_1$ whose inverse image $P''_1$ on~$X''_1$ is
isomorphic to the inverse image of~$P''$. Using the fact that
$\varphi$ is fully faithful, one sees that the object $x_1$ of
$(g^* f_* \Phi)_{Y''_1}$ which corresponds to $P''_1$ is endowed
with a descent datum relative to $Y''_1 \to Y''$, hence
comes from an element $x$ of~$(g^* f_*
\Phi)_{Y''}$. Since the image of~$x$ in $(f'_* \Phi')_{Y''}$ is
$P''$, this proves that $\varphi$ is essentially surjective and
completes the proof.

\begin{example}
\label{XIII.1.4}
Let $f \colon X \to Y$ be a proper morphism. It follows from \cite[VII
2.2.2]{XIII.2} that, for every ind-finite stack $F$ on~$X$, the pair $(F,f)$
is cohomologically proper (relative to $Y$) in dimension
$\leqslant 1$. In particular, for every sheaf of sets
(\resp every sheaf of groups, \resp every ind-finite sheaf of
groups) $F$ on~$X$, $(F,f)$ is cohomologically
% original p. 352
proper in dimension $\leqslant 0$ (\resp in dimension $\leqslant 0$,
\resp in dimension $\leqslant 1$).
\end{example}

\begin{remarks}
\label{XIII.1.5}
a) Let $F$ be a sheaf of groups on~$X$ such that $(F,f)$ is
cohomologically proper relative to~$S$ in dimension
$\leqslant -1$ (\resp $\leqslant 0$). If one regards $F$ as a
sheaf of sets, $(F,f)$ is cohomologically proper
relative to~$S$ in dimension $\leqslant -1$ (\resp $\leqslant
0$), but the converse is false.

For example, let $Y$ be the spectrum of a strictly local discrete
valuation ring with closed point $t$ and generic point
$s$, $f \colon X \to Y$ a nonempty scheme over~$Y$ whose closed
fiber is empty, $F$ a nontrivial constant sheaf of groups
on~$X$, and $P$ a torsor under $F$ such that we have $\H^0(X_s,
\leftexp{P\mkern-4mu}{F}|X_s)=1$. Then $(\leftexp{P\mkern-4mu}{F},f)$ is
cohomologically proper relative to $Y$ in dimension
$\leqslant -1$ when one regards $\leftexp{P\mkern-4mu}{F}$ as a
sheaf of sets. If one regards $\leftexp{P\mkern-4mu}{F}$ as a sheaf
of groups, one has an isomorphism ${\vphantom{F}}^{P^\circ\mkern-7mu}(\leftexp{P\mkern-4mu}{F})
\simeq F$; since the canonical morphism
$$
\H^0(X,F)\to\H^0(X_t,F|X_t)=1
$$
is not injective, this proves that $(\leftexp{P\mkern-4mu}{F},f)$ is not
cohomologically proper relative to $Y$ in dimension
$\leqslant -1$.


b) Suppose $f$ coherent (\ie quasi-compact and
quasi-separated). Let $F$ be a stack on~$X$. For every
geometric point $\overline{y}$ of~$Y'$, we denote by $\overline{Y}$
(\resp $\overline{Y'}$) the strict localization of~$Y$ (\resp $Y'$) at
$\overline{y}$, and we set $\overline{X} = X \times_Y \overline{Y}$,
$\overline{X'} = X' \times_{Y'} \overline{Y'}$, etc. In order that $(F,f)$
be cohomologically proper relative to~$S$ in dimension
$\leqslant -1$ (\resp $\leqslant 0$, \resp $\leqslant 1$), it is necessary and
sufficient that, for every $S$-scheme $S'$ and for every geometric
point $\overline{y}$ of~$Y'$, the canonical functor
$$
\overline{F}(\overline{X})\to\overline{F'}(\overline{X'})
$$
be faithful (\resp fully faithful, \resp an
equivalence).

Indeed,
% original p. 353
if $S'$ is an $S$-scheme, in order that the functor
$$
g^* f_* F \to f'_* F'
$$
be faithful (\resp fully faithful, \resp an
equivalence), it is necessary and sufficient that the same hold for the functor
induced on the fibers at the various geometric points
$\overline{y}'$ of~$\overline{Y'}$ \cite[III 2.1.5.9]{XIII.2}. The assertion
therefore follows from the computation of the geometric fibers of the direct
image of a stack by a coherent morphism \cite[VII 2.1.5]{XIII.2}.


c) Let $F$ be a stack on~$X$. The fact that $(F,f)$ is
cohomologically proper relative to~$S$ in dimension
$\leqslant -1$ (\resp $\leqslant 0$, \resp $\leqslant 1$) is local
on~$Y$ for the \'etale topology.

Let $S'$ be an $S$-scheme, and $F'$ the inverse image of~$F$ on~$X'$
(\cf \eqref{eq:XIII.1.0.1}). If $(F,f)$ is cohomologically proper
relative to~$S$ in dimension $\leqslant 1$, the same holds
for~$(F', f')$. But if $(F,f)$ is cohomologically proper
relative to~$S$ in dimension $\leqslant -1$ (\resp $\leqslant
0$), the same does not necessarily hold for~$(F', f')$.

For example, let $S'$ be a discrete valuation ring, $f' \colon
E_{S'} \to S'$ the affine space over~$S'$, $x$ a closed
point of~$E_{S'}$ lying over the generic point of~$S'$, and
$F'$ the sheaf of sets on~$E_{S'}$ whose restriction to
$E_{S'}-\{x\}$ is the constant sheaf with one element and
whose fiber at a geometric point over~$x$ has two
elements. Then $(F',f')$ is not cohomologically proper
relative to~$S'$ in dimension $\leqslant -1$. Let $S =
S'[Z]$, $f \colon E_S \to S$ the affine space over~$S$, and $T$ a
closed subset of~$X = E_S$ which does not meet the closed set $Z=0$ and
such that $f(T)$ contains the generic point of~$S$. Let
$G$ be the inverse image of~$F'$ on~$X$ and $F$ the sheaf on~$X$ obtained
by extending $G|X-T$ by the empty set. Then $(F,f)$ is
cohomologically proper relative to~$S$ in dimension
$\leqslant -1$, but this is no longer so after the
change of base $S' \to S$ defined by $Z=0$.


d)
% original p. 354
Let $F$ be a stack on~$X$ such that $(F,f)$ is cohomologically
proper relative to $Y$ in dimension $\leqslant -1$ (\resp $\leqslant 0$, \resp $\leqslant 1$). Then, for every
geometric point $\overline{y}$ of~$Y$, the canonical functor
$$
(f_* F)_{\overline{y}} \to F(X_{\overline{y}})
$$
is faithful (\resp fully faithful, \resp an equivalence
of categories).
\end{remarks}

\begin{proposition}
\label{XIII.1.6}
Let $f \colon X \to Y$ and $g \colon Y \to Z$ be two $S$-morphisms, and
$\Phi$ a stack on~$X$.
\begin{enumerate}
\item[1)] Suppose that $(\Phi,f)$ and $(f_* \Phi, g)$ are
cohomologically proper relative to~$S$ in dimension
$\leqslant -1$ (\resp $\leqslant 0$, \resp $\leqslant 1$). Then the
same holds for~$(\Phi, gf)$.
\item[2)] Suppose that $(\Phi, gf)$ is cohomologically proper
relative to~$S$ in dimension $\leqslant -1$ (\resp that $(\Phi,
gf)$ is cohomologically proper relative to~$S$ in dimension
$\leqslant 0$ and $(\Phi,f)$ cohomologically proper relative
to~$S$ in dimension $\leqslant -1$, \resp that $(\Phi, gf)$ is
cohomologically proper relative to~$S$ in dimension
$\leqslant 1$ and $(\Phi, f)$ cohomologically proper relative
to~$S$ in dimension $\leqslant 0$). Then $(f_* \Phi, g)$ is
cohomologically proper relative to~$S$ in dimension
$\leqslant -1$ (\resp in dimension $\leqslant 0$, \resp in dimension
$\leqslant 1$).
\end{enumerate}
\end{proposition}

For every $S$-scheme $S'$, we consider the following diagram,
all of whose squares are cartesian:
\begin{equation*}
\label{eq:XIII.1.6.1}
\tag{\thesubsection.1}
\begin{array}{c}
\xymatrix{ X' \ar[r]^{f'} \ar[d]_h & Y'
\ar[r]^{g'} \ar[d]_k & Z' \ar[r] \ar[d]_{m} & S' \ar[d] \\ X \ar[r]^f
& Y \ar[r]^g & Z \ar[r] & \,S. }
\end{array}
\end{equation*}
The
canonical morphism
$$
m^* (g_* f_* \Phi) \to g'_* f'_* (h^* \Phi)
$$
is identified
% original p. 355
with the composite of the canonical morphisms
$$
m^*(g_* f_* \Phi) \lto{i} g'_*(k^*f_* \Phi) \lto{j} g'_* f'_*(h^* \Phi).
$$
\begin{enumerate}
\item[1)] The hypothesis implies that $i$ and $j$ are
faithful (\resp fully faithful, \resp
equi\-va\-lences); hence the same holds for~$ji$.
\item[2)] The hypothesis implies that $ji$ is faithful
(\resp that $ji$ is
fully
faithful and $j$ faithful, \resp that $ji$ is an equivalence
and $j$ fully faithful); it follows that $i$ is
faithful (\resp fully faithful, \resp an equivalence).
\end{enumerate}

\begin{corollary}
\label{XIII.1.7}
Let $f \colon X \to Y$ and $g \colon Y \to Z$ be two $S$-morphisms,
and let $F$ be a sheaf of groups on~$X$. Suppose that $(F,gf)$
is cohomologically proper relative to~$S$ in dimension
$\leqslant -1$ (\resp that $(F,gf)$ is cohomologically proper
relative to~$S$ in dimension $\leqslant 0$ and that $(F,f)$ is
cohomologically proper relative to~$S$ in dimension
$\leqslant -1$). Then $(f_* F, g)$ is cohomologically proper
relative to~$S$ in dimension $\leqslant -1$ (\resp in dimension
$\leqslant 0$).
\end{corollary}

Let us take up again the notation of \eqref{eq:XIII.1.6.1} and, for every
scheme $Y_1$ \'etale over~$Y$, let us denote by $f_1$, $F_1$ the
respective inverse images of~$f$, $F$ by the morphism $Y_1 \to
Y$. Let $\Phi$ be the stack of torsors under $F$ and $\Psi$ the stack
of torsors under $f_* F$. We have a canonical functor
$$
\varphi \colon \Psi \to f_* \Phi \text{,}
$$
obtained by associating to every scheme $Y_1$ \'etale over~$Y$ and
to every torsor $P$ on~$Y_1$ with group $f_{1 *}F_1$ the torsor
$\tilde{P}$ on~$X_1$ deduced from~$f_1^* P$ by extension of the
structure group $f_1^* f_{1 *}F_1 \to F_1$. The functor
$\varphi$ is fully faithful. Indeed, if $P$ and $Q$ are two
torsors on~$Y_1$ with group $f_{1 *} F_1$, we have a canonical
morphism
$$
\SheafIsom_{f_{1 *}F_1}(P,Q) \to
f_{1*}(\SheafIsom_{F_1}(\tilde{P},\tilde{Q}))
$$
which
% original p. 356
is an isomorphism since this is so locally. We deduce from this
that the canonical morphism
$$
\Isom_{f_{1 *}F_1}(P,Q) \to \Isom_{F_1}(\tilde{P},\tilde{Q})
$$
is an isomorphism, hence that $\varphi$ is fully faithful.

We have a commutative diagram
$$
\xymatrix{ g'_* k^* \Psi \ar[r]^i \ar[d]_{g'_*
k^*(\varphi)} & m^*(g_* \Psi) \ar[d]^{m^*
g_*(\varphi)} \\ g'_* k^*(f_* \Phi) \ar[r]^j &
m^*(g_* f_* \Phi)\text{,} }
$$
where $i$ and $j$ are the change of base morphisms. It
follows from~\Ref{XIII.1.6} 2) that $j$ is faithful
(\resp fully faithful). Since $g'_* k^* (\varphi)$ and
$m^* g_* (\varphi)$ are fully faithful, we deduce
from the above diagram that $i$ is faithful (\resp fully
faithful).

\begin{corollary}
\label{XIII.1.8}
Let $f \colon X \to Y$ be a coherent $S$-morphism, $g \colon Y
\to Z$ a proper $S$-morphism, and $\Phi$ an ind-finite stack on~$X$ \textup{\cite[VII 2.2.1]{XIII.2}}. Suppose that $(\Phi,f)$ is cohomologically proper
relative to~$S$ in dimension $\leqslant -1$ (\resp $\leqslant
0$, \resp $\leqslant 1$); then the same holds for~$(\Phi, gf)$.
\end{corollary}

Since $f$ is coherent, $f_* \Phi$ is an ind-finite stack (SGA~4 IX
1.6~(ii)). The corollary therefore follows from \Ref{XIII.1.6}~1)
and~\Ref{XIII.1.4}.

\begin{corollary}
\label{XIII.1.9}
Let $f \colon X \to Y$ be an integral $S$-morphism, and $g \colon Y \to Z$
an $S$-morphism. If $F$ is a sheaf of sets on~$X$, in order that
$(f_* F, g)$ be cohomologically proper relative to~$S$
in dimension $\leqslant -1$ (\resp $\leqslant 0$), it is necessary and
sufficient that the same hold for~$(F, gf)$. If $F$ is a sheaf of
groups on~$X$, in order that $(f_* F, g)$ be cohomologically
proper relative to~$S$ in dimension $\leqslant -1$
(\resp $\leqslant 0$, \resp $\leqslant 1$), it is necessary and sufficient that
the same hold for~$(F, gf)$.
\end{corollary}

The assertion
% original p. 357
relating to the case of a sheaf of sets follows
from~\Ref{XIII.1.6} and from the fact that $(F,f)$ is cohomologically proper
relative to~$S$ in dimension $\leqslant 0$. Let $F$ be a
sheaf of groups on~$X$ and $\Phi$ the stack of torsors under $F$.
By SGA~4 VIII~5.8, every torsor under $F$ is locally
trivial over~$Y$. It follows that the stack $f_* \Phi$ is
equivalent to the stack of torsors under $f_* F$, the equivalence
being obtained by associating to every scheme $Y_1$ \'etale
over~$Y$ and to every torsor $P$ on~$X_1 = X \times_Y Y_1$ with group
$F|X_1$ the torsor $f_* P$ with group $f_* F | Y_1$. Since $(F,f)$ is
cohomologically proper relative to~$S$ in dimension
$\leqslant 1$, the corollary therefore follows from~\Ref{XIII.1.6}.
